% s-380-p380-m002; original PDF 380
\PoemStart{s-380-p380-m002}{Der Ring.}
\begin{BookVerse}
\VerseLine{0}{Seinen Ring vom braunen Finger}
\VerseLine{1}{Will ihm Colombine streifen.}
\VerseLine{1}{Pierrot schweigt. Die Ruder schleifen.}
\VerseLine{0}{Segel regt kein Kühlebringer.}
\Stanza
\VerseLine{0}{„Mädel sind doch arme Dinger,}
\VerseLine{1}{Männer könnt nach Lüsten schweifen \Dots{3}“}
\VerseLine{0}{Seinen Ring vom braunen Finger}
\VerseLine{1}{Will ihm Colombine streifen.}
\Stanza
\VerseLine{0}{Hände werden heisse Ringer}
\VerseLine{1}{Um den ärmsten aller Reifen.}
\VerseLine{1}{Plötzlich, wie sie dreh’n und greifen, —}
\VerseLine{0}{Schnappt ein Fisch, ein rascher Schlinger,}
\VerseLine{0}{Ihm den Ring vom braunen Finger.}
\end{BookVerse}
\PoemEnd
